\documentclass[11pt]{article}

\usepackage[
  % ----- РАЗМЕРЫ ПОСТЕРА -----
  posterwidth=42cm,
  posterheight=30cm,
  postercolumns=3,
  % ----- СТИЛЬ ШАПОК -----
  % light Светлые (голубые, чёрный текст) - ПО УМОЛЧАНИЮ
  % dark Тёмные (фиолетовые, белый текст) | accent Акцентные (жёлтые)
  boxheaders=dark,
  % ----- СТИЛЬ ТАБЛИЦ / РИСУНКОВ / АЛГОРИТМОВ -----
  boxedtable,        % Таблицы в рамке (иначе минималистичные)
  boxedfigure,       % Рисунки в рамке
  boxedalgorithm,    % Алгоритмы в рамке
  % ----- НАСТРОЙКИ ОТСТУПОВ И РАМОК -----
  % padding=15pt, boxrule=1pt, arc=4mm
  % ----- QR-КОДЫ -----
  % qrdir=qr, qrsize=6em, qrminsize=3.5em, qrgap=1.5em
  % ----- БИБЛИОГРАФИЯ -----
  citingstyle=authoryear,
  bibliostyle=plainnat,
  bibfile=references
]{brainposter}

\usepackage{shortcuts}
\usepackage{pgfplots}
\pgfplotsset{compat=1.18}
\usetikzlibrary{positioning, arrows.meta, calc, fit, backgrounds}

% Палитра графиков берётся из цветов постера, чтобы всё смотрелось единым
\colorlet{plotA}{braincolor}
\colorlet{plotB}{boxcolor!65!black}
\colorlet{plotC}{highlightcolor!75!black}

% Общий стиль для всех графиков — задаём один раз
\pgfplotsset{
  brainplot/.style={
    % 0.98, а не 1.0: pgfplots считает width по осям, подписи торчат наружу
    width=0.98\linewidth, height=4.9cm,
    tick label style={font=\small},
    label style={font=\small},
    legend style={font=\small, draw=gray!40, fill=white, fill opacity=0.85,
                  text opacity=1, row sep=1pt},
    grid=major, grid style={gray!22, thin},
    axis line style={gray!55},
    tick style={gray!55},
    every axis plot/.append style={line width=1pt},
  }
}

\renewcommand{\qrlabelowntg}{Telegram: @jdoe}

\setbrainmeta{
  title={Neural Alignment Techniques for Enhanced Semantic Representations},
  authors={
    Jane Doe\textsuperscript{1}, John Smith\textsuperscript{2}, Alice Johnson\textsuperscript{1}
  },
  affiliations={
    \textsuperscript{1}Moscow Institute of Physics and Technology \\
    \textsuperscript{2}Institute for System Programming, RAS
  }
  % несколько логотипов слева: baselogos {mipt.pdf,isp.pdf,innopolis.pdf}
}

\begin{document}
\begin{mainpart}

% ========================= КОЛОНКА 1 =========================

% \begin{custombox}{Заголовок} ... \end{custombox}
\begin{custombox}{Key Contributions}
\begin{itemize}[leftmargin=1.2em, itemsep=3pt, topsep=2pt]
  \item A \textbf{contrastive alignment} objective that ties text and vision
        encoders into one semantic space.
  \item A closed-form bound on the alignment gap under strong convexity.
  \item \textbf{+4.8\,pp} average Top-1 over the strongest baseline, at
        \textbf{2.3$\times$} lower training cost.
\end{itemize}
\end{custombox}

\section*{Method}

Two frozen encoders are mapped into a shared space by a light projection head;
only the head and the alignment module are trained.

% Рисунки: \begin{figure} ... \captionof{figure}{Подпись} \end{figure}
% Это НЕ плавающий float — блок встаёт ровно там, где написан.
\begin{figure}
% Координаты заданы явно (в см). Так надёжнее, чем `right=... of ...`:
% позиционирование относительно центра узла легко приводит к наложениям.
\begin{tikzpicture}[
  box/.style={rounded corners=2pt, draw=braincolor, fill=braincolor!12,
              minimum height=8mm, minimum width=19mm, align=center,
              font=\small},
  enc/.style={box, fill=braincolor, text=white, draw=braincolor},
  core/.style={box, fill=boxcolor!45, draw=boxcolor!65!black,
               minimum height=14mm},
  % NB: стиль нельзя назвать `out` — так называется встроенный ключ TikZ
  outbox/.style={box, fill=highlightcolor!45, draw=highlightcolor!70!black},
  ar/.style={-{Stealth[length=2mm]}, draw=braincolor!70, line width=0.7pt},
]
  \node[box] (tin)   at (0,    0.85) {Text $x$};
  \node[box] (iin)   at (0,   -0.85) {Image $v$};
  \node[enc] (tenc)  at (2.8,  0.85) {Text\\Encoder};
  \node[enc] (ienc)  at (2.8, -0.85) {Vision\\Encoder};
  \node[core] (align) at (5.6, 0)    {Alignment\\Module};
  \node[outbox] (space) at (8.3, 0)  {Shared\\Space};

  % (a) -- (b) между узлами TikZ сам обрезает по границам — стрелки не залезают
  \draw[ar] (tin)  -- (tenc);
  \draw[ar] (iin)  -- (ienc);
  \draw[ar] (tenc) -- (align);
  \draw[ar] (ienc) -- (align);
  \draw[ar] (align) -- (space);
\end{tikzpicture}
\captionof{figure}{Only the projection head and the alignment module receive
gradients; both encoders stay frozen.}
\label{fig:arch}
\end{figure}

% \begin{assumption} / \begin{theorem} / \begin{definition} / ...
% \makeanchor{имя} ставит якорь, \hyperlink{имя}{текст} — ссылку на него
\begin{assumption}
\makeanchor{assumption:strong_convexity}
The objective \(f\) is $\mu$-strongly convex: for all \( x, y \in X \),
\begin{equation}\label{eq:strong_convexity}
f(y) \geq f(x) + \langle \nabla f(x), y - x \rangle + \frac{\mu}{2} \|y - x\|^2.
\end{equation}
\end{assumption}

Under the \hyperlink{assumption:strong_convexity}{strong convexity assumption},
the alignment gap contracts linearly in the number of steps.

\begin{lemma}[Cauchy--Schwarz]
For all \( u, v \in X \): \( |\langle u, v \rangle| \le \|u\|\,\|v\| \),
with equality iff \( u \) and \( v \) are collinear.
\end{lemma}

\begin{definition}
A map \( T: X \to X \) is \emph{non-expansive} if
\( \|T(x) - T(y)\| \le \|x - y\| \) for all \( x, y \in X \).
\end{definition}

% ========================= КОЛОНКА 2 =========================
% \columnbreak задаёт разрыв колонки явно. Для постера так предсказуемее,
% чем автоматическая балансировка: заголовок не отрывается от своего блока.
% NB: \vfill\null перед \columnbreak НЕ ставим — он растягивает колонку на всю
% высоту, и тогда список литературы под multicols уезжает на вторую страницу.
\columnbreak

\section*{Training}

% \begin{algorithm}{Название}       - ширина колонки
% \begin{algorithm}[ширина]{Название} - узкие центрируются автоматически
% Тело — в algorithmic: [0] без нумерации, [1] каждая строка, [n] каждая n-я
\begin{algorithm}{Projected Alignment Descent}
\begin{algorithmic}[1]
\State \textbf{Input:} point \( x_0 \), step \( \eta \), projection \( \Pi \)
\For{ \( t = 1 \) to \( T \) }
  \State \( g_t \gets \nabla f(x_{t-1}) + \lambda \nabla \mathcal{L}_{\text{align}} \)
  \State \( x_t \gets \Pi(x_{t-1} - \eta\, g_t) \)
\EndFor
\State \Return \( x_T \)
\end{algorithmic}
\end{algorithm}

\begin{figure}
\begin{tikzpicture}
\begin{axis}[
  brainplot,
  xlabel={Training steps, k}, ylabel={Top-1 accuracy, \%},
  xmin=0, xmax=100, ymin=58, ymax=97,
  legend pos=south east,
]
\addplot[plotA, mark=*, mark size=1.1]
  coordinates {(0,61)(10,74)(20,81)(30,85)(40,88)(50,90)(60,92)(70,93.4)(80,94.3)(90,94.9)(100,95.2)};
\addlegendentry{Ours}
\addplot[plotB, mark=square*, mark size=1.1, dashed]
  coordinates {(0,60)(10,70)(20,76)(30,80)(40,83)(50,85)(60,86.5)(70,87.6)(80,88.4)(90,88.9)(100,89.3)};
\addlegendentry{CLIP-style}
\addplot[plotC, mark=triangle*, mark size=1.3, dotted]
  coordinates {(0,59)(10,66)(20,71)(30,74)(40,77)(50,79)(60,80.5)(70,81.6)(80,82.3)(90,82.8)(100,83.1)};
\addlegendentry{Frozen baseline}
\end{axis}
\end{tikzpicture}
\captionof{figure}{Validation accuracy on CIFAR-10. Our objective converges
faster and to a higher plateau.}
\label{fig:curves}
\end{figure}

\begin{theorem}[Linear contraction]
Let Assumption~1 hold and let \( \eta \le 1/L \). Then
\( \|x_T - x^\star\|^2 \le (1 - \eta\mu)^T \|x_0 - x^\star\|^2 \).
\end{theorem}

\begin{figure}
\begin{tikzpicture}
\begin{axis}[
  brainplot,
  height=4.3cm,
  ybar, bar width=7pt,
  ylabel={Top-1 accuracy, \%},
  symbolic x coords={CIFAR-10, SVHN, STL-10},
  xtick=data, ymin=70, ymax=100,
  legend style={at={(0.5,1.03)}, anchor=south, legend columns=3},
  nodes near coords={}, % подписи над столбцами можно включить здесь
]
\addplot[fill=plotA, draw=plotA] coordinates {(CIFAR-10,95.2)(SVHN,97.1)(STL-10,93.6)};
\addlegendentry{Ours}
\addplot[fill=plotB, draw=plotB] coordinates {(CIFAR-10,89.3)(SVHN,94.1)(STL-10,86.2)};
\addlegendentry{CLIP-style}
\addplot[fill=plotC, draw=plotC] coordinates {(CIFAR-10,83.1)(SVHN,90.4)(STL-10,79.8)};
\addlegendentry{Frozen}
\end{axis}
\end{tikzpicture}
\captionof{figure}{Consistent gains across three datasets.}
\label{fig:bars}
\end{figure}

% ========================= КОЛОНКА 3 =========================
\columnbreak

\section*{Results}

% Таблицы: \begin{table} ... \captionof{table}{Подпись} \end{table}
% Выделение ячеек: \highlightcell \highlightcellA \highlightcellB \highlightcellC
% Выделение строк:  \highlightrow  \highlightrowA  \highlightrowB  \highlightrowC
\begin{table}
\centering
\begin{tabular}{l c c c}
\toprule
Model & CIFAR-10 & SVHN & Params \\
\midrule
MLP \rule{0pt}{2.3ex} & 84.5\% & \highlightcell 91.2\% & 1.2M \\
CNN-3 &\highlightcellC 89.3\% & 94.1\% & 0.8M \\
\highlightrowB
ResNet18 & 92.7\% & 95.6\% & 11.7M \\
ResNet152 & 95.2\% & 97.1\% & 36.5M \\
\highlightrow
ViT-Tiny & 93.1\% & 96.8\% & 5.6M \\
\bottomrule
\end{tabular}
\captionof{table}{Accuracy and parameter count across datasets.}
\label{table:results}
\end{table}

\begin{figure}
% Матрица ошибок нарисована обычным TikZ: значение задаёт и число, и насыщенность
\begin{tikzpicture}[x=9.5mm, y=9.5mm, font=\small]
  \foreach \c [count=\i from 0] in {cat, dog, car, ship} {
    \node[anchor=east] at (0, 3.5-\i) {\c};
    \node[rotate=40, anchor=east] at (\i+0.5, -0.15) {\c};
  }
  % NB: рамку каждой клетки рисуем вместе с заливкой (\filldraw), а не общей
  % сеткой \draw...grid: у grid шаг по умолчанию step=1cm — абсолютный, он не
  % учитывает x=/y= выше, и сетка расходится с клетками.
  \foreach \x/\y/\v in {
    0/3/91, 1/3/6, 2/3/2, 3/3/1,
    0/2/7, 1/2/89, 2/2/3, 3/2/1,
    0/1/1, 1/1/2, 2/1/94, 3/1/3,
    0/0/1, 1/0/1, 2/0/5, 3/0/93}
  {
    \filldraw[fill=braincolor!\v, draw=gray!45, line width=0.3pt]
      (\x, \y) rectangle ++(1, 1);
    \node[text=\ifnum\v>55 white\else black\fi] at (\x+0.5, \y+0.5) {\v};
  }
  \node[anchor=north, font=\small] at (2,-1.05) {predicted};
\end{tikzpicture}
\captionof{figure}{Confusion matrix (\%), rows are ground truth.}
\label{fig:cm}
\end{figure}

% QR-коды: \addqr{файл=подпись, ...} — размер и перенос считаются сами.
% Готовый набор лаборатории — \qrcodes. Файлы лежат в qr/ (см. qr/README.md).
\begin{qrbox}{Scan for code \& contacts}
\addqr{x-qr.png=Lab on X, lab-tg.png=\qrlabelowntg}
\end{qrbox}

\nocite{*}  % включает все работы из .bib файла
\end{mainpart}
\end{document}
